% Page budget: 1.55 pages | Owns: augmentation topology, additional states, encoder, and closed-loop objective.
% WRITING GUIDE
% Purpose: Define the final augmentation architecture and explain how it is trained inside the known feedback loop.
% Start here: Current implementation decisions D-129, D-140, D-141, D-180, D-220, D-233, and D-237 in docs/decisions.md.
% Supporting material: Quinten's 18 September outline feedback, the Meeting-audit synthesis, Thesis-writeup/Documentation/TRAINING-DESIGN.md, and docs/writeup figures.
% Mandatory papers: Read Hoekstra's 2025 LFR augmentation paper, the 2026 first-principles augmentation paper, the encoder-initialization paper, and Kessels' Chapter 5 before writing this section.
% Research-plan anchor: Preserve Aspect 2's dynamic parallel augmentation objective; document the departure from open-loop training and earlier state routing.
% Current decisions: Separate physical and additional states, keep the controller outside the model, score the complete training window without burn-in, and use one network for the physical correction and added-state update.
% Choices to justify: Final reported routing, added-state dimension, ANN width and depth, zero initialization, encoder history, closed-loop objective, rollout length, full-window scoring, and validation horizon. Use the configuration of the reported thesis run, not a superseded routing experiment.
% Horizon check: Treat encoder history, scored training window, joint-estimation horizon, and free-run validation horizon separately; relate candidates to relevant stable modes and test sensitivity to slower dynamics and free-integrator accumulation.
% Implementation audit: Read scripts/gantry/gantry_interconnect_dynamic.py, gantry_dynamic/{model,config,controller,data,training}.py, and model_augmentation/fit_systems/{interconnect,closed_loop,pre_encoder}.py.
% Reference check: Verify sources for parallel LFR augmentation, encoder initialization, closed-loop state extension, and any claimed architectural precedent against the original paper or thesis.
% Cautions: Earlier routing plans are superseded; distinguish the truth-only hidden absorber from learned states; closed-loop fit alone does not prove autonomous plant recovery.
% Figures: Absorber schematic, author's updated architecture, and thesis-owned common-reference loops. No residual panel or horizon figure.
% Section is finished when: Every signal has a defined frame and timing, the RK4 boundary is explicit, and the described architecture matches the final code.
%
% SOURCE MAP
% Hoekstra et al. 2025 EJC: dynamic-parallel topology, additional states, encoder-based truncated simulation.
% Hoekstra et al. 2026 Automatica submission: extended LFR formulation, baseline-preserving model/encoder initialization, joint identification context.
% Hoekstra et al. 2026 encoder paper: reconstructability-based W^b initialization; it does not identify W^a or test dynamic augmentation.
% Kessels 2025 thesis: direct closed-loop rollout, controller equations, and reconstruction of the machine controller state for each window.
% Drenth 2025 thesis: LPV-specific augmentation precedent; not the unrestricted MLP structure used here.
% This project: gantry specialization, residual-loop implementation, exact routing, RK4 boundary, and verification.
\section{Dynamic LPV-LFR Augmentation}\label{sec:augmentation}

The baseline of Section~\ref{sec:lfr} describes the rigid-body dynamics of
the gantry, but not every effect of the machine.
Garc\'ia-Herreros et al.\ neglect base resonance, cross-arm vibration on
the flexible plates, motor detent forces, and friction variation along the
stroke~\cite[Sec.~2.4]{garcia2013model}.
Our baseline \eqref{eq:eom} also leaves out the Coulomb friction that their
model contains.
This section augments the baseline with a learned dynamic component that is
intended to represent such effects.
The physical parameters $\theta_{\mathrm{base}}$, the network parameters
$\theta_{\mathrm{aug}}$ and the encoder parameters $\eta$ are estimated
jointly from records measured under feedback.
Section~\ref{sec:aug_structure} defines the augmented model,
Section~\ref{sec:aug_encoder} the encoder that estimates the initial state of
each training window, Section~\ref{sec:aug_closed_loop} the closed-loop
objective, and Section~\ref{sec:aug_init} the initialisation.

\subsection{Augmented model}\label{sec:aug_structure}
An omitted vibration mode requires additional model order.
Refitting $\theta_{\mathrm{base}}$ cannot add this order, and neither can a
static correction of the state update, because both only change the dynamics
of the existing states.
Additional states can add the missing order~\cite[Sec.~2]{hoekstra2025lfr}.
We connect the learned part in parallel to the baseline, because the
orthogonality construction of Section~\ref{sec:obc} acts on an additive
correction.
We therefore identify the dynamic-parallel
model~\cite[Table~1]{hoekstra2026lfr}
\begin{equation}
\begin{aligned}
 \tilde x_{k+1}
 &=f_{\mathrm{base}}(\theta_{\mathrm{base}},\tilde x_k,\hat u_k)
   +f_{\mathrm{aug}}(\theta_{\mathrm{aug}},\hat x_k,\hat u_k),\\
 \bar x_{k+1}
 &=g_{\mathrm{aug}}(\theta_{\mathrm{aug}},\hat x_k,\hat u_k),\\
 \hat y_k&=C_n\tilde x_k ,
\end{aligned}
\label{eq:aug_transition}
\end{equation}
where $\tilde x=\col(q,\dot q)\in\R^6$ is the physical state of
Section~\ref{sec:lfr}, $\bar x\in\R^{n_{x_a}}$ collects the additional
states, $\hat x=\col(\tilde x,\bar x)$, $\hat y_k$ is the predicted
position, $C_n$ is the output map, and $\hat u_k$ is the actuator force
applied to the model.
Section~\ref{sec:aug_closed_loop} derives $\hat u_k$ from the recorded
actuator forces $u_k=\uact$ and the measured positions $y_k=\qact$.
The baseline transition $f_{\mathrm{base}}$ keeps its physical parameters.
The learned functions $f_{\mathrm{aug}}$ and $g_{\mathrm{aug}}$ act in
parallel to it and receive the full state $\hat x_k$.

The baseline transition $f_{\mathrm{base}}$ is one step of the classical
fourth-order Runge-Kutta (RK4) method over the sample period $T_s$, applied
to the state equation of \eqref{eq:eom}.
The input is held over the step, as the actuator force is held between
samples.
Each RK4 stage evaluates $M(Y)$ at its own state.
As a result, the scheduling variable $Y$ moves within the step, and
$f_{\mathrm{base}}$ remains self-scheduled.
Drenth formulates the augmentation of such a self-scheduled LPV-LFR
baseline, including additional states~\cite[Sec.~5.2]{drenth2025thesis}.
Inputs, outputs and states are normalised to zero mean and unit standard
deviation with statistics of the training
records~\cite[Sec.~5.3]{hoekstra2026lfr}.
The RK4 step is applied in these normalised coordinates, which changes the
coordinates but not the dynamics.

The two learned functions are the outputs of one multilayer perceptron (MLP)
with tanh hidden layers,
\begin{equation}
 \begin{bmatrix}
  f_{\mathrm{aug}}(\theta_{\mathrm{aug}},\hat x_k,\hat u_k)\\
  g_{\mathrm{aug}}(\theta_{\mathrm{aug}},\hat x_k,\hat u_k)
 \end{bmatrix}
 =\phi_{\mathrm{aug}}\bigl(\theta_{\mathrm{aug}},\col(\hat x_k,\hat u_k)\bigr),
 \label{eq:aug_mlp}
\end{equation}
where $\theta_{\mathrm{aug}}$ collects the weights and biases of the MLP
$\phi_{\mathrm{aug}}$.
The first six outputs form $f_{\mathrm{aug}}$ and correct every
physical-state row, so the correction reaches $X$ and $Y$ as well as
$\Theta$.
A correction of the $\Theta$ rows alone is not sufficient, because the
dynamics missing from the baseline couple into the $X$ and $Y$ responses.
The last $n_{x_a}$ outputs form $g_{\mathrm{aug}}$ and give the next
additional state.
Because the network output is added to the result of the RK4 step,
$f_{\mathrm{aug}}$ is a discrete-time correction.
The baseline, in contrast, integrates the continuous-time physics.
The width and depth of the MLP are given in Section~\ref{sec:setup}.

\begin{figure}[t]
 \centering
 \resizebox{\columnwidth}{!}{\input{figures/tex/augmentation_structure}}
 \caption{Augmented model. The physics step $f_{\mathrm{base}}$ and the
 network $\phi_{\mathrm{aug}}$ act in parallel on $\hat x_k$;
 $\phi_{\mathrm{aug}}$ supplies $f_{\mathrm{aug}}$ and $g_{\mathrm{aug}}$,
 and the output map is not learned ($h_{\mathrm{aug}}=0$). The encoder
 $\psi$ sets $\hat x$ at each window start. During training, $u_k$ is the
 closed-loop input $\hat u_k$.}
 \label{fig:aug_structure}
\end{figure}

The output map is not learned (Fig.~\ref{fig:aug_structure}).
It is the kinematic map $C_n=\begin{bmatrix}P^\top & 0\end{bmatrix}$ of
\eqref{eq:Ptransform} in normalised coordinates, so the measured positions
are read from the physical state alone.
Within one step, no learned function feeds the input of another, so the
interconnection graph is acyclic.
Because $M(Y)$ is invertible over the whole stroke, $f_{\mathrm{base}}$ is
smooth, and so are the tanh layers.
Together with the acyclic graph, this makes the model well
posed~\cite[Thm.~9]{hoekstra2026lfr}.
The additional states have no assigned physical meaning and reach the output
only through $f_{\mathrm{aug}}$.
In particular, they are not the states of the absorber that the simulated
truth contains (Fig.~\ref{fig:gantry_absorber}).

\subsection{Encoder and initial state}\label{sec:aug_encoder}
Each training window simulates \eqref{eq:aug_transition} from an initial
state, but only the positions are measured.
Following Beintema et al.~\cite{beintema2023subnet}, an encoder estimates
this state from the measured output and input history,
\begin{equation}
 \hat x_{\tau|\tau}
 =\psi_\eta(\xi_\tau),\qquad
 \xi_\tau=\col\bigl(y_{\tau-n_a:\tau},\,u_{\tau-n_b:\tau}\bigr),
 \label{eq:aug_encoder}
\end{equation}
where $\tau$ is the window start, $y_{\tau-n_a:\tau}$ stacks the outputs
$y_{\tau-n_a},\dots,y_\tau$, and $u_{\tau-n_b:\tau}$ stacks the inputs in the
same way.
Both histories thus include the current sample.
The encoder estimates all states, including the measured positions.
The history lengths are given in Section~\ref{sec:setup}.
The encoder is a linear map with a nonlinear
correction~\cite[Eq.~(8)]{hoekstra2026encoder},
\begin{equation}
 \psi_\eta(\xi_\tau)
 =\begin{bmatrix}W^b\\ W^a\end{bmatrix}\xi_\tau+\tilde\psi(\xi_\tau),
 \label{eq:aug_enc_lin}
\end{equation}
where $W^b$ and $W^a$ give the physical and additional states and
$\tilde\psi$ is an MLP.
The linear part can be initialised from the baseline, as
Section~\ref{sec:aug_init} shows.
The nonlinear correction can then learn the part of the state reconstruction
that the linear map misses.

\subsection{Closed-loop rollout and objective}\label{sec:aug_closed_loop}
We fit the closed-loop response of the model, because the records are
measured under feedback and the model is used under feedback.
In that use, the model predicts the machine's response to a given reference,
feedforward and known controller.
This departs from the open-loop simulation driven by the recorded input
in~\cite[Eq.~(22)]{hoekstra2026lfr}.
The baseline gives a second reason.
$X$ and $Y$ have no stiffness in \eqref{eq:damping_stiffness_matrices}, so
the linearised baseline has two eigenvalues at the origin for every $Y$.
In an open-loop rollout, a position error caused by a velocity error
therefore persists.
Inside the loop, the controller acts on such errors as it does on the
machine, provided the model in feedback with the controller is stable.

Kessels trains an augmented model with the known feedback controller inside
each truncated rollout and propagates gradients through
it~\cite[Eqs.~(5.12), (5.13c), (5.13d)]{kessels2025ai}.
We adopt this rollout.
The controller stays outside the learned model, so that it can be replaced
in evaluation~\cite[Sec.~5.3.2.3]{kessels2025ai}.
The recorded loop is driven by a reference $r$ and a feedforward
$u^{\mathrm{ff}}$ through a linear controller with state $s$.
The model loop receives the same $r$, $u^{\mathrm{ff}}$ and controller.
It differs only in its output $\hat y$ (Fig.~\ref{fig:aug_closed_loop}):
\begin{equation}
\begin{aligned}
 s_{k+1}&=A_k s_k+B_k(r_k-y_k), &
 u_k&=C_k s_k+D_k(r_k-y_k)+u^{\mathrm{ff}}_k,\\
 \hat s_{k+1}&=A_k \hat s_k+B_k(r_k-\hat y_k), &
 \hat u_k&=C_k \hat s_k+D_k(r_k-\hat y_k)+u^{\mathrm{ff}}_k ,
\end{aligned}
\label{eq:aug_direct_controller}
\end{equation}
where $A_k$, $B_k$, $C_k$ and $D_k$ are the known controller matrices and
$\hat s$ is the state of the model's controller.

\begin{figure}[t]
 \centering
 \includegraphics[width=\columnwidth]{closed_loop_same_reference.pdf}
 \caption{Recorded plant loop (top) and augmented-model loop (bottom),
 driven by one shared reference $r_k$ and feedforward $u_k^{\mathrm{ff}}$,
 with the same feedback controller. The model output is evaluated
 before the input advances its state to the next sample.}
 \label{fig:aug_closed_loop}
\end{figure}

Subtracting the recorded loop from the model loop removes $r$ and
$u^{\mathrm{ff}}$,
\begin{equation}
\begin{aligned}
 x^c_{k+1}&=A_k x^c_k+B_k(y_k-\hat y_k),\qquad x^c_\tau=0,\\
 \hat u_k&=u_k+C_k x^c_k+D_k(y_k-\hat y_k),
\end{aligned}
\label{eq:aug_controller}
\end{equation}
where $x^c=\hat s-s$ is the difference of the controller states.
The model thus receives the recorded input plus the controller's response to
its own output error.
This form needs only the recorded $u$ and $y$, not $r$ and
$u^{\mathrm{ff}}$.
It is exact under three conditions: both loops use the same linear
controller and feedforward, the controller is not scheduled on the model
state, and neither loop saturates.

The condition $x^c_\tau=0$ starts the model's controller in the state of the
recorded controller, $\hat s_\tau=s_\tau$, without reconstructing that
state.
As a result, model error before $\tau$ is not carried into the window.
Kessels instead reconstructs the controller state from the measured output,
the reference and the controller~\cite[Remark~5.4]{kessels2025ai}.
Within the window, the augmented model has no feedthrough, so each sample is
evaluated in a fixed order.
The output $\hat y_k$ follows from the state, then $\hat u_k$ from
\eqref{eq:aug_controller}, and then both states advance.

Over window starts $\mathcal T$ and window length $n_f$, the parameters
$\vartheta=(\theta_{\mathrm{base}},\theta_{\mathrm{aug}},\eta)$ minimise the
output error
\begin{equation}
 V(\vartheta)=\frac{1}{|\mathcal T|\,n_f}
 \sum_{\tau\in\mathcal T}\sum_{\ell=0}^{n_f-1}
 \norm{y_{\tau+\ell}-\hat y_{\tau+\ell|\tau}}_2^2
 \label{eq:aug_loss}
\end{equation}
subject to \eqref{eq:aug_transition}, \eqref{eq:aug_encoder} and
\eqref{eq:aug_controller}, where $\hat y_{\tau+\ell|\tau}$ is the
closed-loop prediction started from $\hat x_{\tau|\tau}$.
This is the truncated objective of~\cite[Eq.~(22)]{hoekstra2026lfr}, with
each window simulated in closed loop as in~\cite[Eq.~(5.12)]{kessels2025ai}.
The physical parameters, the network and the encoder are estimated jointly,
while the controller is fixed.
As in~\cite[Eq.~(22)]{hoekstra2026lfr}, every sample of the window is
scored.

Three horizons enter this objective and its use, and they are set
separately.
The encoder history $n_a$, $n_b$ sets how much past data informs the initial
state.
The window length $n_f$ sets how far a model error propagates before it is
scored.
It must span the memory of the closed loop, through which a force error of
the model reaches the output.
Validation, and with it checkpoint selection, uses closed-loop simulation
over the complete validation records, which are much longer than one window.
Section~\ref{sec:setup} gives the values and the basis of each.

A small closed-loop error does not show that the model is accurate in open
loop.
The controller acts on the model's output error as it acts on the
machine's, so feedback can hide model error.
The closed-loop objective therefore does not establish on its own that the
plant dynamics are recovered.
Section~\ref{sec:results} tests this separately from the closed-loop fit.

\subsection{Initialisation}\label{sec:aug_init}
Random initialisation of the learned parts can make the augmented model
unstable from the start~\cite[Sec.~5.4]{hoekstra2026lfr}.
Training therefore starts from the baseline.
To this end, we initialise the output layer of the network at zero,
\begin{equation}
 W_L=0,\quad b_L=0
 \quad\Longrightarrow\quad
 f_{\mathrm{aug}}=0,\quad g_{\mathrm{aug}}=0 ,
 \label{eq:aug_zero_init}
\end{equation}
where $W_L$ and $b_L$ are the weight and bias of the output layer of
$\phi_{\mathrm{aug}}$.
The augmented model then reproduces the baseline from the same physical
initial state, which is the initial behaviour required
in~\cite[Eq.~(29)]{hoekstra2026lfr}.

The encoder must also start from baseline behaviour: at initialisation, it
should return the baseline state~\cite[Sec.~5.4.2]{hoekstra2026lfr}.
We therefore initialise the physical rows $W^b$ from a linearisation of the
baseline.
We linearise \eqref{eq:eom} at its nominal parameters about rest at $Y=0$,
which gives
\begin{equation}
\begin{aligned}
 A_c(Y)&=\begin{bmatrix}0 & I\\ -M(Y)^{-1}K & -M(Y)^{-1}C\end{bmatrix},\qquad
 B_c(Y)=\begin{bmatrix}0\\ M(Y)^{-1}P\end{bmatrix},\\
 C_d&=\begin{bmatrix}P^\top & 0\end{bmatrix},
\end{aligned}
\label{eq:aug_lin}
\end{equation}
where $A_c(Y)$ and $B_c(Y)$ are the state and input matrices at frozen $Y$
and $C_d$ is the output matrix.
The pair $(A_d,B_d)$ is the zero-order-hold discretisation of
$(A_c(0),B_c(0))$ over $T_s$, because the input is held between samples.

For this linear model, $W^b$ is the reconstructability
map~\cite[Eqs.~(16), (17)]{hoekstra2026encoder}
\begin{equation}
 W^b=\begin{bmatrix}
  A_d^nO_n^{+} & \;r_n-A_d^nO_n^{+}T_n
 \end{bmatrix},
 \label{eq:aug_wb}
\end{equation}
where $n=n_a=n_b$ is the encoder history, $O_n$ is the observability matrix
of $(A_d,C_d)$ over $n+1$ samples, $T_n$ is the Toeplitz matrix of
input-to-output responses, and $r_n$ is the input-to-state map over the same
samples.
The two blocks act on the output and input parts of $\xi_\tau$.
The measured positions make the linear model observable, so $O_n$ has full
column rank and $O_n^{+}$ is a left inverse.
All matrices are normalised in the same way as the signals.
The linearisation only sets the starting value of $W^b$, because the encoder
is trained jointly with the model.

The rows $W^a$ have no baseline counterpart and are drawn randomly from a
Kaiming-uniform distribution.
This departs from the Xavier initialisation
of~\cite[Sec.~5.4.2]{hoekstra2026lfr}.
The nonlinear correction $\tilde\psi$ starts at zero.
By \eqref{eq:aug_zero_init}, the encoded additional states do not affect the
output at initialisation.
They become active as the output layer of the network is trained.

From this start, $f_{\mathrm{aug}}$ can also learn changes that a different
$\theta_{\mathrm{base}}$ would produce.
The split between the two components is therefore not
unique~\cite[Sec.~5.2]{hoekstra2026lfr}.
Section~\ref{sec:obc} constrains this split.
The number of additional states is given in Section~\ref{sec:setup}.
